% Generated draft: requires mathematical and source review.
\documentclass[11pt]{article}
\usepackage[margin=25mm]{geometry}
\usepackage{fontspec}
\setmainfont{Latin Modern Roman}
\usepackage{amsmath,amssymb,mathtools,unicode-math}
\setmathfont{Latin Modern Math}
\usepackage{enumitem,needspace,xurl,hyperref}
\setlength{\parindent}{0pt}
\setlength{\parskip}{5pt}
\setlength{\emergencystretch}{3em}
\newcommand{\N}{\mathbb{N}}
\newcommand{\Z}{\mathbb{Z}}
\newcommand{\Q}{\mathbb{Q}}
\newcommand{\R}{\mathbb{R}}
\newcommand{\C}{\mathbb{C}}
\newcommand{\F}{\mathbb{F}}
\newcommand{\A}{\mathbb{A}}
\newcommand{\PP}{\mathbb{P}}
\newcommand{\E}{\mathbb{E}}
\newcommand{\eps}{\varepsilon}
\newcommand{\dd}{\,\mathrm{d}}
\begin{document}
\section*{UnsolvedMath ID: 2858}
\textbf{Source problem code:} catalogue entry
\subsection{Definitions and mathematical statement}
The imported task has two targets: compute instanton Floer homology of closed $3$-manifolds and compute Donaldson invariants of closed $4$-manifolds. Here closed means compact without boundary. The usual gauge-theoretic versions also require smooth structures, orientations and auxiliary choices, none of which the short statement spells out.

Instanton Floer theory is based on the Chern--Simons functional on connections modulo gauge; in appropriate versions its differential counts anti-self-dual trajectories on a cylinder. Its precise group depends on the chosen version, bundle/admissibility data, handling of reducible connections, coefficient ring and grading convention. It is not a uniquely specified group for every closed $3$-manifold without such a convention.

Donaldson invariants are gauge-theoretic invariants formed using anti-self-dual moduli spaces and prescribed cohomological insertions. Their formulation requires the relevant bundle/topological data, orientation information and, in low-$b_2^+$ settings, chamber choices or restrictions. The source does not specify these or a version covering every named manifold.

\textbf{Source-faithful algorithmic objective.} Construct terminating procedures that, from a finite encoding of each admissible manifold and the data defining the desired invariant, return an exact algebraic description of the chosen instanton Floer group or the requested Donaldson invariant. For example, once the invariant and its admissible domain have been fixed, the intended computability assertions would have the form
\[
 \mathcal A_I(\operatorname{enc}(Y),\mathcal d)\longmapsto I(Y;\mathcal d),
 \qquad
 \mathcal A_D(\operatorname{enc}(X),\mathcal e,z)\longmapsto D_{X,\mathcal e}(z),
\]
with termination for every valid input. Here $Y$ has dimension three, $X$ has dimension four, $\mathcal d,\mathcal e$ stand for the required invariant-specific data, and $z$ stands for insertions if an evaluation is requested. These are specification placeholders, not an invented choice of Floer version or Donaldson theory.

A rigorous algorithm statement must fix the input encoding, such as appropriate surgery/handle diagrams, specify the smooth structure in dimension four, and state the output representation including coefficients and grading. It must also decide whether Donaldson computation means any specified evaluation or a finite description of an entire invariant. No finite-output description of an otherwise infinite family is assumed. The source does not provide these choices, so the broad request cannot yet be turned into one fully checkable computational theorem.

The literature notes and triage dated 2026-08-17 report algorithms for Heegaard Floer invariants of knots, $3$-manifolds and cobordisms, including mixed invariants of closed $4$-manifolds. They cite Manolescu--Ozsváth--Sarkar, Sarkar--Wang, Lipshitz--Manolescu--Wang and Manolescu--Ozsváth--Thurston. These concern Heegaard Floer theories, not an established algorithm for the requested instanton/Donaldson targets. The background says Problem 3.58 and Witten's Conjecture should relate these to Donaldson theory for manifolds of simple type; that conditional and restricted comparison is not promoted to a general algorithm for all closed $4$-manifolds.

The dated review reports no verified algorithm for the requested targets and calls for direct or comparison-based constructions. This is a literature-status report, not a proof that such algorithms do not exist. The registry label is open and the triage label partially solved. No algorithm, PDF or comparison theorem was independently checked, and no present-day assessment is claimed.
\subsection{Short English statement}
Develop exact terminating algorithms for instanton Floer homology in dimension three and Donaldson invariants in dimension four. The input and output conventions and auxiliary gauge-theory data need specification. Existing Heegaard Floer algorithms and conjectural simple-type comparisons are reported as related progress, not as solutions to these targets.
\subsection{Sources}
\begin{itemize}
\item UnsolvedMath local registry, data/unsolvedmath/problems.json, permanent ID 2858, source code catalogue entry, ``Kirby Problem 3.60,'' including the two algorithm targets and literature triage dated 2026-08-17.
\item \emph{K3: A New Problem List in Low-Dimensional Topology}, Problem 3.60 (2026), supplied link \url{https://aimath.org/pastworkshops/kirbylistrep.pdf}. The review cites Heegaard Floer algorithms and outstanding targets; the PDF was not opened here.
\item [MOS09] Ciprian Manolescu, Peter Ozsváth and Sucharit Sarkar, ``A combinatorial description of knot Floer homology,'' \emph{Ann. of Math. (2)} 169(2) (2009), 633--660, DOI 10.4007/annals.2009.169.633.
\item [SW10] Sucharit Sarkar and Jiajun Wang, ``An algorithm for computing some Heegaard Floer homologies,'' \emph{Ann. of Math. (2)} 171(2) (2010), 1213--1236. The DOI is supplied as 10.4007/annals.2010.171.12s13 and is retained without independent validation or guessed correction.
\item [MOT25] Ciprian Manolescu, Peter Ozsváth and Dylan Thurston, ``Grid diagrams and Heegaard Floer invariants,'' \emph{Ann. of Math. (2)} 201(1) (2025), 1--78, DOI 10.4007/annals.2025.201.1.1.
\item [LMW08] Robert Lipshitz, Ciprian Manolescu and Jiajun Wang, ``Combinatorial cobordism maps in hat Heegaard Floer theory,'' \emph{Duke Math. J.} 145(2) (2008), 207--247, DOI 10.1215/00127094-2008-050.
\item Ulam AI and the UnsolvedMath Contributors, \emph{UnsolvedMath: A Curated Collection of Open Mathematics Problems}. Dataset curation and original metadata are supplied under CC BY 4.0; referenced source material retains its own terms. Problem 3.58 and the section introduction containing Witten's Conjecture are not reproduced, and no external reference was independently consulted for this draft.
\end{itemize}
\end{document}
